% FILE: 06_contribution_to_thesis.tex
% PROJECT: sources-papers
% ROLE: academic-paper-classification-and-evidence-grounding

\section{Contribution to the Dissertation}
\label{sec:sources-papers-contribution}

\subsection{Contribution to Prediction vs.~Coordination}
\label{sec:sources-papers-contribution-prediction}

The dissertation's central argument is that the dominant paradigm of enterprise
AI --- prediction --- is insufficient for process-critical systems, and that
\emph{coordination} grounded in verifiable process evidence is required. The
paper corpus substantiates this argument with two concrete demonstrations from
the academic literature.

First, the predictive process monitoring papers catalogued in
\texttt{PAPER\_TO\_TYPE\_LAW.md} --- Tax et al.\ (2017) LSTM-based prefix
prediction, Di Francescomarino et al.\ (2017) survey, Polato et al.\ (2018)
time-aware prediction --- share a structural limitation: their output is
\texttt{PredictionTarget}, a single predicted value, not a \emph{process
evidence chain}. The \texttt{PrefixLength} and \texttt{PredictionHorizon} types
required to make these predictions typed and bounded are documented as MISSING
from wasm4pm-compat. A system built on prediction alone cannot produce the
\texttt{Admitted<OcelLog, Ocel20>} $\to$ discovery $\to$ conformance $\to$
board claim chain that the dissertation's architecture requires.

Second, the adversarial canon analysis in \texttt{adversarial-canon.md}
demonstrates that prediction is not merely insufficient but actively dangerous
in adversarial environments. Alpha Miner, Heuristics Miner, and Inductive
Miner --- the canonical discovery algorithms --- produce incorrect models under
Sybil attack because they predict process structure from event frequency
patterns without verifying event provenance. Coordination via BLAKE3-receipted
event streams is the required correction, not a more sophisticated predictor.

\subsection{Contribution to Process Evidence}
\label{sec:sources-papers-contribution-evidence}

The paper corpus makes four specific contributions to the dissertation's theory
of process evidence:

\begin{enumerate}
  \item \textbf{The Paper-to-Law mapping algorithm} $\mathcal{M}: I_P \to
    (\mathcal{T}_{\text{Rust}}, \mathcal{W})$ is the first formal account, in
    this research program, of how academic invariants become compiler-enforced
    guarantees. The mapping is not interpretive; it is mechanical: every
    structural invariant that can be expressed without dynamic graph traversal
    becomes a named Rust type, and every violation of that invariant becomes a
    compile-time error with a diagnostic citing the originating paper.

  \item \textbf{The 5-valued dynamic valuation lattice}
    $\mathcal{V} = \{\bot, \text{PossiblySatisfied}, \text{Satisfied},
    \text{Violated}, \top\}$ from \texttt{declare-satisfaction-lattice.md}
    provides the mathematical foundation for streaming process evidence
    production. Evidence is not a binary verdict but a monotonically advancing
    lattice element: once \texttt{Satisfied} or \texttt{Violated} is reached,
    the state is permanent and immutable. This property is essential for receipt
    chain integrity --- a receipt cannot be retroactively invalidated by a later
    event.

  \item \textbf{The \texttt{Between01<NUM, DEN>} conformance metric bound}
    makes fitness, precision, and F1 measurements typed at the compile level.
    The distinction between a measurement (typed, bounded, receipted) and an
    assertion (untyped floating-point) is the core distinction between process
    evidence and process narrative.

  \item \textbf{The PARTIAL coverage map} --- 22 named missing types, 14 named
    missing output shapes --- is itself a contribution to process evidence
    theory. It demonstrates that the obligation surface is enumerable, that
    gaps can be named precisely, and that the boundary between ``covered'' and
    ``not covered'' is a compile-time property, not an opinion.
\end{enumerate}

\subsection{Contribution to Manufacturing Architecture}
\label{sec:sources-papers-contribution-manufacturing}

The CodeManufactory thesis argues that process artifacts --- models, receipts,
conformance scores, board claims --- must be manufactured, not asserted. The
paper corpus contributes to this argument by defining the manufacturing inputs.

The four core lifecycle algorithms in \texttt{core\_lifecycle\_algorithms.md}
define the manufacturing pipeline for process evidence:

\begin{enumerate}
  \item \textbf{Paper-to-Law}: Converts academic invariants into type
    obligations (the specification layer of the manufacturing pipeline).
  \item \textbf{PM4Py Oracle Mapping}: Maps PM4Py algorithm capabilities to
    wasm4pm gap status, establishing the oracle reference for the manufacturing
    engine's capability surface.
  \item \textbf{Admission Gate}: Enforces the \texttt{Admitted<T, W>}
    invariant, preventing raw event laundering from entering the manufacturing
    pipeline.
  \item \textbf{Receipt-Bearing Execution}: Cryptographically binds each
    manufacturing output to the evidence chain, making the provenance of every
    artifact non-forgeable.
\end{enumerate}

The graduation boundary taxonomy --- \texttt{NeedsDiscovery},
\texttt{NeedsConformanceExecution}, \texttt{NeedsObjectCentricQueryExecution},
\texttt{NeedsReceipts}, \texttt{RebuildingProcessMiningLocally} --- maps
directly to manufacturing stage classifications: each graduation reason names
the manufacturing capability that must be present before a paper's algorithms
can be manufactured as receipted execution artifacts.

\subsection{Contribution to Capital-Flow Infrastructure}
\label{sec:sources-papers-contribution-capital}

The dissertation's long-horizon claim is that verifiable process evidence is a
prerequisite for capital-flow infrastructure --- settlement systems, DAOs,
regulatory compliance chains --- that requires machine-readable proof of
process conformance, not human-readable assertions about process conformance.

The paper corpus contributes to this claim by establishing the board-admissible
claim taxonomy in \texttt{PAPER\_TO\_BOARD\_CLAIM.md}. A board-admissible
claim, as defined there, must satisfy three conditions: grounded in a specific
paper with a specific formal result; verifiable evidence path in
\texttt{wasm4pm-compat} (named Rust type, compile-fail fixture, or graduation
boundary); not falsifiable by pointing to a missing type that contradicts the
claim.

This three-condition definition is the process-evidence analogue of the legal
definition of admissible evidence. It establishes that the output of the
manufacturing pipeline --- conformance receipts, fitness scores, alignment
results --- can function as machine-readable evidence in capital-flow contexts
where human attestation is insufficient. The 5 Tier-1 structural completeness
claims currently documented as board-admissible represent the minimal viable
evidence surface for a process-intelligence-backed capital-flow claim; the 5
claims explicitly marked NOT board-admissible due to missing types document
the remaining manufacturing obligation.
